%=====================================================================
\section{Methods}
\label{sec:methods}
Every analysis follows the same template: hypothesis from the literature $\to$ metric
$\to$ randomized reference model. \Cref{tab:hypotheses} lists the seven hypotheses,
the methods (M1--M4) that test them and the datasets where each method is meaningful
(19 method--dataset pairs). The empty cells are deliberate: CIT is a year-resolution
DAG without fine-grained events (no M2a, M2b, M3) and without community
reorganisation (no M4), while PS and CM have no growth process (no M2c).

\begin{table*}[t]
\centering
\small
\setlength{\tabcolsep}{4pt}
\begin{tabular*}{\textwidth}{@{\extracolsep{\fill}}cl p{8cm} l l l ccc@{}}
\toprule
Part & ID & Hypothesis & Source & Method & Null & PS & CM & CIT \\
\midrule
I   & H1 & Aggregation and window choice change structural conclusions
           & \citep{krings2012effects,ribeiro2013quantifying} & M1 & -- & \checkmark & \checkmark & \checkmark \\
\midrule
II  & H2 & Activity is bursty and has memory
           & \citep{goh2008burstiness} & M2a & P1 & \checkmark & \checkmark & \\
    & H3 & Ties are repeated and persistent; the network turns over
           & \citep{saramaki2014persistence} & M2b & P1 & \checkmark & \checkmark & \\
\midrule
III & H4 & Growth follows preferential attachment with ageing; networks densify
           & \citep{jeong2003measuring,wang2013quantifying,leskovec2005graphs} & M2c-1--4 & P3 & & & \checkmark \\
    & H5 & Communities reorganise over time; the static partition hides this
           & \citep{greene2010tracking,mucha2010community} & M4 & P1 & \checkmark & \checkmark & \\
\midrule
IV  & H6 & Time ordering restricts reachability and changes centrality rankings
           & \citep{holme2005network,pan2011path,kim2012temporal} & M3a, M3b & P1, P2 & \checkmark & \checkmark & \\
    & H7 & Burstiness and temporal correlations slow down spreading
           & \citep{karsai2011small} & M3c & P1, P2 & \checkmark & \checkmark & \\
\bottomrule
\end{tabular*}
\caption{Hypotheses tested in this work, with methods, null models
(\cref{sec:null}) and applicable datasets.}
\label{tab:hypotheses}
\end{table*}

\subsection{Notation}
\label{sec:notation}
Let $\mathcal{E}=\{(u,v,t)\}$ be the event sequence. The static aggregate $G$
contains a link $(u,v)$ if at least one event occurs between $u$ and $v$. A
snapshot $G_{\Delta t}(k)$ contains the events in $[k\Delta t,(k+1)\Delta t)$, and the
cumulative network $G(\le t)$ contains all events up to $t$.

\subsection{Aggregation and time windows (M1)}
\label{sec:m-aggregation}
For each window length $\Delta t$ on a grid (PS: 5\,min to 1\,day; CM: 1\,h to 1\,month)
we compute density, mean degree, clustering coefficient and giant-component size of the
snapshots and of the aggregate. For CIT, yearly and cumulative networks are compared.
\emph{Window sensitivity} is the relative change of a metric across the grid.
Rankings are handled separately in M3b. Aggregation can introduce paths that
cannot be traversed in time: if $c$--$d$ occurs before $a$--$b$ and $b$--$c$,
the static path $a$--$b$--$c$--$d$ is not time-respecting from $a$.




\subsection{Activity patterns (M2a, M2b)}
\label{sec:m-activity}

\textbf{M2a: Burstiness and memory.}
We analyse inter-event times for each interacting node pair
in the Primary School (PS) and CollegeMsg (CM) networks.
PS contacts are recorded every 20 seconds. Consecutive
records for the same unordered pair, separated by exactly
20 seconds, are merged into a single contact episode.
For CM, each retained directed message is treated as
an independent event.

For each edge, inter-event intervals are calculated
from consecutive event timestamps. We pool these
intervals across edges and calculate burstiness as
\begin{equation}
B = \frac{\sigma-\mu}{\sigma+\mu},
\end{equation}
where $\mu$ and $\sigma$ are the mean and population
standard deviation of the pooled intervals.

Memory is defined as the Pearson correlation between
consecutive inter-event intervals within each edge,
pooled across all valid adjacent interval pairs
\citep{goh2008burstiness}.
Both $B$ and $M$ are pooled statistics rather than
averages of per-edge coefficients.

\textbf{M2b: Repetition, lifetime and turnover.}
We quantify temporal link persistence using three metrics.

The repetition rate is the fraction of distinct edges
with at least two independent contact episodes (PS)
or two messages (CM).

Edge lifetime is defined as the difference between
the last and first event timestamps of each edge.
Edges observed only once have zero lifetime.

Link turnover is measured using the Jaccard similarity
between edge sets in consecutive time windows:
\begin{equation}
J(E_k,E_{k+1}) =
\frac{|E_k \cap E_{k+1}|}
{|E_k \cup E_{k+1}|}.
\end{equation}
We average this quantity over adjacent window pairs
with nonempty unions.

The default window size is one hour for PS and
seven days for CM. We additionally evaluate
15, 30, 60 and 120-minute windows for PS,
and 1, 3, 7 and 14-day windows for CM.

For PS, edge lifetime and window Jaccard are computed
from the original event stream. We also compare raw
20-second records with merged contact episodes
to assess the sensitivity of H2 to event definition.


\subsection{Structural evolution (M2c, M4)}
\label{sec:m-evolution}
% Owner: B (expanded M2c-1 through M2c-4 and P3).
\paragraph{M1 for citations.}
Annual and cumulative citation networks share the same node set at year $t$, namely all
selected papers published by $t$. Directed density is $E/[N(N-1)]$ and mean total degree
is $2E/N$; clustering, component counts and LCC fractions use the weak undirected
projection, including isolates. Trailing windows of 1, 2, 5 and 10 years provide the
window grid. Window sensitivity, $(\max-\min)/\mathrm{mean}$ over the grid, is computed at
each end year and averaged over end years 1990--2020, which constitute the primary
analysis interval; earlier papers enter as history.

\paragraph{M2c-1: preferential attachment.}
Following the degree-class approach of \citet{jeong2003measuring}, in-degrees are frozen
at the start of each citing year, and for each prior-degree class $k$ the number of
eligible papers $n_{tk}$ and of new citations $h_{tk}$ are counted. The exponent $\alpha$
is estimated by conditional likelihood with target probabilities proportional to
$n_{tk}(k+1)^\alpha$, conditioning on the total number of citations in each year. The
offset admits citations to previously uncited papers; offsets of 0.5 and 2 serve as
sensitivity checks. Confidence intervals use 400 bootstrap resamples of citing-year blocks
and are descriptive. Neither a positive $\alpha$ nor its departure from an age-only
reference identifies a causal mechanism.

\paragraph{M2c-2: ageing.}
For ages 1--60, the target probability in citing year $t$ is proportional to
$n_{ta}f(a)$, where $n_{ta}$ is the number of available papers of age $a$. Lognormal,
exponential, power-law and uniform kernels $f$ are compared by AIC on identical
observations and support. The comparison is repeated by publication cohort (1970--1989,
1990--1999, 2000--2009), accounting for each cohort's exposure. The estimate describes
marginal ageing, which may also reflect degree, fitness and venue coverage.

\paragraph{M2c-3: densification and diameter.}
The densification exponent $a$ is obtained from $\log E(t)=a\log N(t)+b$ over 1990--2020
and over split periods and an earlier cutoff. Year-resampling intervals are descriptive
because cumulative observations are dependent. The effective diameter is the 90th
percentile of shortest-path distances within the largest weak component, estimated from
128 uniformly sampled source nodes per year, with 64 and 256 sources as sensitivity checks
\citep{leskovec2005graphs}.

\paragraph{M2c-4: citation dynamics and prediction.}
The Wang--Song--Barab\'asi (WSB) model \citep{wang2013quantifying},
\begin{equation}
 c(t)=m\left[\exp\!\left\{\lambda\Phi\left(\frac{\log t-\mu}{\sigma}\right)\right\}-1\right],
\end{equation}
is fitted with $m$ fixed to the mean within-subgraph reference count of papers published
up to 1989. Bounded multistart least squares minimizes the squared error of $\log(1+c)$;
parameter bounds and boundary flags are recorded. Descriptive fits use a random sample of
at most 500 papers published in 1990--2005 with at least 10 citations by age 15, fitted
over ages 1--15; the resulting collapse is in-sample evidence only. For prediction,
papers are selected using only their first five years (at least five citations by age
five; random sample of at most 500), and WSB is compared with a no-ageing preferential
attachment baseline, $c(t)=m[\exp(rt)-1]$, and with linear extrapolation, $c(5)t/5$. The
primary task predicts cumulative citations at age 15, which is observed for every selected
paper; a secondary task predicts the 2020 total with a variable horizon. Neither target
represents lifetime citations. This fixed-horizon design differs from the original study,
which trains on 5 or 10 years and evaluates at age 30 \citep{wang2013quantifying}.
Performance is measured by mean absolute error (MAE), median absolute error and root mean
squared logarithmic error (RMSLE), with paired bootstrap differences (1000 resamples). Fits
at a parameter bound remain in the primary evaluation, and an interior-fit subset serves
as a sensitivity check; a constrained optimum is not regarded as a convergence failure.

\paragraph{P3: citation-target randomization.}
Twenty realizations per model preserve each paper's publication year and its number of
distinct references, which are drawn without replacement from strictly older papers. P3a
selects targets uniformly. P3b weights targets by a nonparametric age kernel fitted by
conditional likelihood to the observed data and removes degree preference. Because P3b is
calibrated on the data, its agreement with the observed ageing is not independent
confirmation. Null means and sample standard deviations are reported; empirical tail
probabilities have a resolution of $1/21$. Both models fix $N(t)$ and $E(t)$ and thus
preserve the densification exponent exactly; they test diameter and target preference but
not densification itself.


\textbf{M4} detects communities in every snapshot (Leiden or Louvain), matches them
across snapshots \citep{greene2010tracking} and reports modularity and the number of
communities over time, the normalised mutual information (NMI) between consecutive
snapshots, the NMI between the dynamic and the static partition, and for PS the NMI
with the class labels. The resolution parameter is fixed and its effect is checked.

\subsection{Paths, centrality and spreading (M3)}
\label{sec:m-function}
\textbf{M3a} computes earliest-arrival time-respecting paths and the temporal reachability
ratio $R$ (fraction of node pairs connected by a time-respecting path, relative to the
static aggregate). \textbf{M3b} compares static and temporal closeness rankings by
Kendall's $\tau$. \textbf{M3c} runs an event-driven SIR process on the event sequence with
at least 100 random seeds; SI with certain transmission is equivalent to M3a, so only SIR
is run.


\subsection{Randomized reference models}
\label{sec:null}
No single temporal counterpart of the configuration model exists, so one switches off
selected correlations with different models and compares the outcomes
\citep{holme2012temporal,gauvin2022randomized}.

For PS and CM, we consider two complementary types of randomization,
corresponding to \emph{randomly permuted times} (RP) and
\emph{randomized edges} (RE) in the naming of
\citet{holme2012temporal}.

\textbf{P1 (RP): Timestamp permutation.}
P1 randomly permutes timestamps among events while keeping their
endpoints unchanged. It preserves the static graph, the number of
contacts of every pair, and the global timestamp distribution,
while disrupting the original temporal correlations.
For H2 and H3, we generate 20 realizations and report
the null mean and standard deviation.

\textbf{P2: Structural randomization.}
We distinguish two variants of edge randomization.

\textbf{P2a (RE, original proposal).}
The original method rewires the static graph in a
degree-preserving way while each pair's contact sequence
moves with its link. It aims to destroy the wiring diagram
while preserving per-link temporal activity patterns
and the overall event rate.
This variant is retained as the original methodological
proposal but was not evaluated in the Member C experiments.

\textbf{P2b (constrained endpoint swaps, implemented).}
Our implementation randomly swaps endpoints between event pairs
while keeping timestamps fixed. For undirected PS, endpoint
orientations are randomized before swapping; for directed CM,
sender and receiver roles are preserved. Swaps creating self-loops
or duplicate edge-time events are rejected.

P2b preserves the total event count, timestamp multiset,
and node event participation counts (separate sender and
receiver counts for CM). Unlike P2a, it does not necessarily
preserve static simple-graph degrees or individual edge
contact sequences.

We attempt five swaps per event and validate the structural
invariants using a fixed random seed of 42.
These checks do not establish uniform sampling or sufficient
mixing of the randomization process.

Comparing temporal and structural randomization helps
distinguish the effects of event timing from network topology.
Combined models (RE+RP) and fully random times (RT)
remain possible extensions.


\subsection{Measuring the effect of time}
\label{sec:time-effect}
Raw metrics cannot be compared across a contact network, a message network and a
citation network, so we define one normalised measure per hypothesis:
H1 window sensitivity (relative change of a metric over the $\Delta t$ grid);
H2 differences in pooled burstiness and memory
relative to P1,
$B_{\mathrm{real}}-\overline{B}_{\mathrm{P1}}$
and
$M_{\mathrm{real}}-\overline{M}_{\mathrm{P1}}$;
H3 differences in adjacent-window link Jaccard
relative to P1,
$J_{\mathrm{real}}-\overline{J}_{\mathrm{P1}}$;
H4 exponents $\alpha$ and $a$;
H5 mean consecutive-snapshot NMI and dynamic-vs-static NMI;
H6 reachability ratio $R$ and Kendall $\tau$(static, temporal);
H7 slowdown $T_{\mathrm{real}}/T_{\mathrm{null}}$, where $T$ is the time to the epidemic peak.

\subsection{Implementation}
All code is in Python. We use NetworkX for static graph analysis, Raphtory
\citep{steer2024raphtory} for time-respecting paths, tnetwork for dynamic communities,
and Leiden/Louvain for per-snapshot partitions. Null models, the attachment-kernel
estimator and the SIR process are custom code, unit-tested on hand-built examples
whose answers are known (reachability, SIR, preserved quantities of P1 and P2). Random
seeds are fixed and every figure is generated by one script.
